% =============================================================================
\section{Triangulation Related Modules}
\label{sec:triangulationModules}
% =============================================================================

% -----------------------------------------------------------------------------
\subsection{2D Triangulation Bindings}
\label{ssec:modules:tri2}
% -----------------------------------------------------------------------------
A triangulation of a set of points $\calP$ in \Rtwo is a partition of the convex hull of $\calP$ into triangles whose vertices are the points of $\calP$. Together with the unbounded face having the convex hull boundary as its frontier, the triangulation forms a partition of \Rtwo. Any two facets (2-face) are either disjoint or share a common edge (1-face) or vertex (0-face). A triangulation can be described as a simplicial complex. The binding module \iiDTriangulationsPackage{} consists of bindings of types provided by the \cgal{} packages \iiDTriangulationsPackage{} and \iiDPeriodicTriangulationsPackage{}. Table~\ref{tab:tri2} lists the \cmake{} flags associated with the \iiDTriangulationsPackage{} module.

%%%%%%%% 2D Triangulation Module Flags
\begin{table}[!htp]
  \caption{\iiDTriangulationsPackage{} module flags.}
  \centering\begin{tikzpicture}
  \matrix (first) [tableModuleFlags,
    column 1/.style={nodes={text width=16.0em}},
    column 3/.style={nodes={text width=5.7em}},
    column 4/.style={nodes={text width=21.6em}},
    row 6/.append style={nodes={minimum height=2.8em}}
  ] {
    Name & Type & Default & Description\\
    \cmakeLstinline{TRI2\_NAME} & String & \myLstinline{plain} &
      The 2D triangulation type\\
    \cmakeLstinline{TRI2\_VERTEX\_WITH\_INFO}      & Boolean & \myLstinline{false} & Determines whether the vertex type is extented\\
    \cmakeLstinline{TRI2\_FACE\_WITH\_INFO}        & Boolean & \myLstinline{false} & Determines whether the face type is extented\\
    \cmakeLstinline{TRI2\_INTERSECTION\_TAG\_NAME} & String  & \myLstinline{ncirc} & The intersection tag\\
    \cmakeLstinline{TRI2\_HIERARCHY}               & Boolean & \myLstinline{false} & Determines whether to generate the binding for a hierarchy triangulation\\
  };
  \end{tikzpicture}
  \label{tab:tri2}
\end{table}

The \iiDTriangulationsPackage{} package~\cite{cgal:y-t2-21a} provides several class templates, instances of which can be used to represent a variety of 2D triangulations. In particular, the package provides the following class templates:
\begin{compactenum}
\item \cppLstinline{Triangulation\_2<Traits, Tds>},
\item \cppLstinline{Regular\_triangulation\_2<Traits, Tds>},
\item \cppLstinline{Delaunay\_triangulation\_2<Traits, Tds>},
\item \cppLstinline{Constrained\_triangulation\_2<Traits, Tds, Itag>},
\item \cppLstinline{Constrained\_Delaunay\_triangulation\_2<Traits, Tds, Itag>}, and
\item \cppLstinline{Triangulation\_hierarchy\_2<Triangulation\_2>}
\end{compactenum}
Instance of the \cppLstinline{Delaunay\_triangulation\_2<>} and \cppLstinline{Regular\_triangulation\_2<>} class templates can be used to represent Delaunay and regular triangulation, respectively. In a regular triangulation points have an associated weight, and some points can be hidden and do not result in vertices in the triangulation. The class template \cppLstinline{Triangulation\_hierarchy\_2} enables fast point location queries. When instantiated its template parameter must be substituted with an instance of any other triangulation class template.

The \iiDPeriodicTriangulationsPackage{} package~\cite{cgal:k-pt2-13-21a} supports triangulations of sets of points in the two-dimensional flat torus~\cite{ct-cpt-09}. This package provides the class templates
\begin{compactenum}
\item \cppLstinline{Periodic\_2\_triangulation\_2<Traits, Tds>},
\item \cppLstinline{Periodic\_2\_Delaunay\_triangulation\_2<Traits, Tds>}
\item \cppLstinline{Periodic\_2\_triangulation\_hierarchy\_2<PeriodicTriangulation>}
\end{compactenum}

The \cmake{} flag \cmakeLstinline{TRI2\_NAME} specifies the particular type of triangulation, bindings for which should be generated; see Table~\ref{tab:tri2:names}.

\begin{table}[!htp]
  \caption{2D triangulation name options}
  \centering\begin{tikzpicture}
  \matrix (first) [tableName,
    column 1/.style={nodes={text width=11.5em}},
    column 2/.style={nodes={text width=37.5em}}
  ] {
    \cmakeLstinline{TRI2\_NAME} & Type\\
    \myLstinline{plain} & \cppLstinline{Triangulation\_2<Traits, Tds>}\\
    \myLstinline{regular} & \cppLstinline{Regular\_triangulation\_2<Traits, Tds>}\\
    \myLstinline{delaunay} & \cppLstinline{Delaunay\_triangulation\_2<Traits, Tds>}\\
    \myLstinline{constrained} & \cppLstinline{Constrained\_triangulation\_2<Traits, Tds, Itag>}\\
    \myLstinline{constrainedDelaunay} & \cppLstinline{Constrained\_Delaunay\_triangulation\_2<Traits, Tds, Itag>}\\
    \myLstinline{periodicPlain} & \cppLstinline{Periodic\_2\_triangulation\_2<Traits, Tds>}\\
    \myLstinline{periodicDelaunay} & \cppLstinline{Periodic\_2\_Delaunay\_triangulation\_2<Traits, Tds>}\\
  };
  \end{tikzpicture}
  \label{tab:tri2:names}
\end{table}

%% \cppLstinline{TriangulationTraits\_2}
%% \cppLstinline{RegularTriangulationTraits\_2}
%% \cppLstinline{DelaunayTriangulationTraits\_2}
%% \cppLstinline{ConstrainedTriangulationTraits\_2}
%% \cppLstinline{ConstrainedDelaunayTriangulationTraits_2}
%% \cppLstinline{AlphaShapeTraits\_2}
%% \cppLstinline{WeightedAlphaShapeTraits\_2}
%% \cppLstinline{Periodic\_2TriangulationTraits\_2}
%% \cppLstinline{Periodic\_2DelaunayTriangulationTraits\_2}
%% \cppLstinline{Periodic\_3AlphaShapeTraits\_2}
When any template above is instantiated the template parameter \cppLstinline{Traits} must be substituted with a model of a suitable geometric traits concept; this model is referred to as the geometric traits class; it provides the type of points to use as well as elementary operations on them. The type of traits used for the generated bindings is determined based on the selection of the triangulation type as explained below. It is conveniently defined in C++ as \cppLstinline{Dt::Geom\_traits}, where \cppLstinline{Dt} is a triangulation instance. This type is exposed as a Python attribute with the same name under \pyLstinline{Triangulation\_2} (which in turn is nested under the Python namespace \pyLstinline{Tri2}). Figure~\ref{tab:tri2:concept-hierarchy} depicts the 2D triangulation traits concept hierarchy. Any kernel instance is a model of any non-periodic traits concept; thus, when one of the templates
\begin{compactenum}
\item \cppLstinline{Triangulation\_2<Traits, Tds>},
\item \cppLstinline{Regular\_triangulation\_2<Traits, Tds>},
\item \cppLstinline{Delaunay\_triangulation\_2<Traits, Tds>},
\item \cppLstinline{Constrained\_triangulation\_2<Traits, Tds, Itag>}, or
\item \cppLstinline{Constrained\_Delaunay\_triangulation\_2<Traits, Tds, Itag>}
\end{compactenum}
is instantiated, the \cppLstinline{Traits} parameter is substituted
with the \cppLstinline{Kernel} type (the selected kernel; see
Section~\ref{ssec:modules:kernel}). When one of the templates
\begin{compactenum}
\item \cppLstinline{Periodic\_2\_triangulation\_2<Traits, Tds>} or
\item \cppLstinline{Periodic\_2\_Delaunay\_triangulation\_2<Traits, Tds>}
\end{compactenum}
is instantiated, the \cppLstinline{Traits} parameter is substituted with \cppLstinline{Periodic\_2\_triangulation\_traits\_2<Kernel>} or \cppLstinline{Periodic\_2\_Delaunay\_triangulation\_traits_2<Kernel>}, respectively.
%
Observe that when a Delaunay triangulation instance is used to define an alpha shape type (see Section~\ref{ssec:modules:as2}), the traits parameter must be substituted with a traits class that models the \cppLstinline{AlphaShapeTraits\_2} concept. Similarly, when a regular triangulation instance is used to define an fixed alpha shape type (see Section~\ref{ssec:modules:as2}, the traits parameter must be substituted with a traits class that models the \cppLstinline{WeightedAlphaShapeTraits\_2} concept. Also in these cases the selected kernel serves as the traits.

\begin{figure}[!htp]
  \tikzset{
    my node style/.style={font=\scriptsize,minimum size=6mm,align=center,
      top color=white, bottom color=blue!25, draw=blue!75, drop shadow,
      rectangle, rounded corners, very thick
    }
  }
  \forestset{
    my tree style/.style={
      for tree={my node style,
        parent anchor=south,
        child anchor=north,
        edge={->,>=stealth,black,thick},
        edge path={% |-| style branches from parent to child
          \noexpand\path [draw, \forestoption{edge}] (!u.parent anchor)--(.child anchor)\forestoption{edge label};
        },
        if n children=3{% if a node has 3 children
          for children={% align the middle child with its parent
            if n=2{calign with current}{}
          }
        }{},
      }
    }
  }
  \centering
  \begin{forest}
    my tree style
    [ {\cppLstinline{TriTr2}}
      [ {\cppLstinline{RegularTriTr2}}
        [ {\cppLstinline{WeightedAlphaShapeTr2}} ]
      ]
      [ {\cppLstinline{Pr2TriTr2}},name=p ]
      [ {\cppLstinline{DelaunayTriTr2}},name=d
        [ {\cppLstinline{Pr2DelaunayTriTr2}},name=pd
	  [ {\cppLstinline{Pr2AlphaShapeTr2}} ]
        ]
        [ {\cppLstinline{AlphaShapeTr2}} ]
      ]
      [ {\cppLstinline{ConstrainedTriTr2}}
        [ {\cppLstinline{ConstrainedDelaunayTriTr2}},name=cd ]
      ]
    ]
    \draw[->,>=stealth] (d.parent anchor)--(cd.child anchor);
    \draw[->,>=stealth] (p.parent anchor)--(pd.child anchor);
  \end{forest}
  \caption{The 2D triangulation traits concept hierarchy.
    \cppLstinline{Triangulation}, \cppLstinline{Periodic\_2} and
    \cppLstinline{Traits\_2} are abbreviated as \cppLstinline{Tri},
    \cppLstinline{Pr2}, and \cppLstinline{Tr2},
    respectively.}
  \label{tab:tri2:concept-hierarchy}
\end{figure}

The type that substitutes the \cppLstinline{Tds} parameter models the concept
\cppLstinline{TriangulationDataStructure\_2}. An object of this type stores the combinatorial structure of the triangulation; it is an instance of the class template \cppLstinline{Triangulation\_data\_structure\_2<VertexBase, FaceBase>}.

The types that substitute the \cppLstinline{VertexBase} and \cppLstinline{FaceBase} template parameters when the template \cppLstinline{Triangulation\_data\_structure\_2<VertexBase, FaceBase>} is instantiated represent the type of the vertex and the type of the face of the triangulation, respectively; they must model the concepts \cppLstinline{TriangulationDSVertexBase\_2} and \cppLstinline{TriangulationDSFaceBase\_2}, respectively. If the binding is generated for a periodic triangulation, these parameters
must be substituted with types that also model the concepts \cppLstinline{Periodic\_2TriangulationDSVertexBase\_2} and \cppLstinline{Periodic\_2TriangulationDSFaceBase\_2}, respectively. If the triangulation is used to define an alpha shape type, these parameters must be substituted with types that also model the concepts \cppLstinline{AlphaShapeVertex\_2} and \cppLstinline{AlphaShapeFace\_2}, respectively; see Section~\ref{ssec:modules:as2}. Finally, if the binding is generated for a hierarchy triangulation, e.g., an instance of the template parameter \cppLstinline{Triangulation\_hierarchy\_2<Triangulation\_2>}, the \cppLstinline{VertexBase} template parameter must be substituted with a model of the concept \cppLstinline{TriangulationHierarchyVertexBase\_2}. (Observe that there is no special requirements on the type that substitutes the \cppLstinline{FaceBase} parameter in this case.) The final vertex type is selected accordingly as follows.

%% 2D Triangulation Vertex Selection
The type \cppLstinline{Vb} is defined as the vertex base type.  If bindings is generated for a periodic triangulation, the type \cppLstinline{Vb} is defined as \cppLstinline{Periodic\_2\_triangulation\_vertex\_base\_2<Traits>}; otherwise, if bindings is generated for an instance of \cppLstinline{Regular\_triangulation\_vertex\_base\_2<Kernel>}, the
type \cppLstinline{Vb} is defined as \cppLstinline{Regular\_triangulation\_vertex\_base\_2<Traits>}; otherwise, the type \cppLstinline{Vb} is defined as \cppLstinline{Triangulation\_vertex\_base\_2<Traits>}.
%
If the \cmake{} flag \cmakeLstinline{TRI2\_VERTEX\_WITH\_INFO} is set to \myLstinline{true}, the type \cppLstinline{Vbi} is defined as \cppLstinline{Triangulation\_vertex\_base\_with\_info\_2<Data, Kernel, Vb>}, where \cppLstinline{Data} is defined as the generic Python object \cppLstinline{bp::object}; see Section~\ref{sec:code:ae}; otherwise the type \cppLstinline{Vbi} is defined as \cppLstinline{Vb}.
%
If the \cmake{} flag \cmakeLstinline{TRI2\_HIERARCHY} is set to \myLstinline{true}, the type \cppLstinline{Vbih} is defined as \cppLstinline{Triangulation\_hierarchy\_vertex\_base\_2<Vbi>}; otherwise the type \cppLstinline{Vbih} is defined as \cppLstinline{Vbi}.
%
Finally, if the triangulation is used to define an alpha shape type, the final type \cppLstinline{V} (that substitutes the \cppLstinline{VertexBase} parameter) is defined as
\cppLstinline{Alpha\_shape\_vertex\_base\_2<Traits, Vbih, ExactComparison>}, where \cppLstinline{ExactComparison} is substituted with either \cppLstinline{true} or \cppLstinline{false} based on the setting of the \cmake{} flag \cmakeLstinline{AS2\_EXACT\_COMPARISON}; see Table~\ref{tab:as2}. Otherwise, \cppLstinline{V} is defined as \cppLstinline{Vbih}.
%
The final vertex type is conveniently defined in C++ as \cppLstinline{Dt::Vertex}, where \cppLstinline{Dt} is the triangulation instance. This type is exposed as a Python attribute
with the same name under \pyLstinline{Triangulation_2}.

The final face type is determined in a similar fashion. The only difference between the selections of the vertex and face types is that the setting of the \cmake{} flag \cmakeLstinline{TRI2\_HIERARCHY} does not have an effect on the face type selection. The final face type is conveniently defined in C++ as \cppLstinline{Dt::Face}, where \cppLstinline{Dt} is the triangulation instance. This type is also exposed as a Python attribute with the same name under \pyLstinline{Triangulation_2}.

The \cmake{} Boolean flag \cmakeLstinline{TRI2\_HIERARCHY} indicates whether the type of triangulation, bindings for which should be generated, is an instance of one of the triangulation hierarchy templates, namely, \cppLstinline{Triangulation\_hierarchy\_2<Triangulation\_2>} and
\cppLstinline{Periodic_triangulation\_hierarchy\_2<PeriodicTriangulation\_2>}. The template parameter in both cases is substituted with the triangulation selected via the \cmake{} flag
\cmakeLstinline{TRI2\_NAME}, which also determines whether to use the non-periodic or periodic version above.

% -----------------------------------------------------------------------------
\subsection{2D Alpha Shape Bindings}
\label{ssec:modules:as2}
% -----------------------------------------------------------------------------
The alpha shape (a.k.a.\ alpha complex) of a set of points is one of several notions of a shape formed by the set. Given a set of points sampled in a 2D body, an alpha shape is demarcated by a frontier, which is a linear approximation of the original boundary of the body. A two-dimensional alpha shape object maintains an underlying triangulation of a set of input points in the plane. There are two distinguished versions of alpha shapes as follows. Basic alpha shapes are based on the Delaunay triangulation and weighted alpha shapes are based on its generalization, the regular triangulation, where the euclidean distance is replaced by the power to weighted points. The package \iiDAlphaShapesPackage{}~\cite{cgal:d-as2-21a} provides the class template \cppLstinline{Alpha\_shape\_2<Dt, ExactAlphaComparisonTag>}. In a 2D alpha shape object represented by an instance of this class template each $k$-simplex of the underlying triangulation is associated with an interval that specifies for which values of $\alpha$ the $k$-simplex belongs to the alpha shape. Table~\ref{tab:as2} lists the \cmake{} flags associated with the \iiDAlphaShapesPackage{} module.

%%%%%%%% 2D Alpha Shape Module flags
\begin{table}[!htp]
  \caption{\iiiDAlphaShapesPackage{} module flags}
  \centering\begin{tikzpicture}
  \matrix (first) [tableModuleFlags,
    column 1/.style={nodes={text width=16em}},
    column 3/.style={nodes={text width=5.7em}},
    column 4/.style={nodes={text width=21.6em}}
  ] {
    Name & Type &  Default & Description\\
    \myLstinline{AS2\_EXACT\_COMPARISON} & Boolean & \myLstinline{false} &
      Determines whether to apply exact comparisons\\
  };
 \end{tikzpicture}
 \label{tab:as2}
\end{table}

% -----------------------------------------------------------------------------
\subsection{3D Triangulation Bindings}
\label{ssec:modules:tri3}
% -----------------------------------------------------------------------------
A triangulation of a set of points $\calP$ in \Rthree is a partition of the convex hull of $\calP$ into tetrahedra whose vertices are the points of $\calP$. Similar to the two-dimensional triangulation, together with the unbounded cell having the convex hull boundary as its frontier, the triangulation forms a partition of \Rthree. Any two cells (3-face) are either disjoint or share a common facet (2-face), edge (1-face) or vertex (0-face).  The binding module \iiiDTriangulationsPackage{} consists of bindings of types provided by the \cgal{} packages \iiiDTriangulationsPackage{} and \iiiDPeriodicTriangulationsPackage{}. Table~\ref{tab:tri3} lists the \cmake{} flags associated with the \iiiDTriangulationsPackage{} module.

%%%%%%%% 3D Triangulation Module Flags
\begin{table}[!htp]
  \caption{\iiiDTriangulationsPackage{} module flags. For the default
  value of the traits name see Table~\ref{tab:tri3:traits-default}.}
  \centering\begin{tikzpicture}
  \matrix (first) [tableModuleFlags] {
    Name & Type & Default & Description\\
    \cmakeLstinline{TRI3\_NAME}                   & String & \myLstinline{plain}    & The 3D triangulation type\\
    \cmakeLstinline{TRI3\_CONCURRENCY\_NAME}      & String & \myLstinline{sequential} & The concurrency method\\
    \cmakeLstinline{TRI3\_LOCATION\_POLICY\_NAME} & String & \myLstinline{compact}    & The location policy\\
  };
  \end{tikzpicture}
  \label{tab:tri3}
\end{table}

The \iiiDTriangulationsPackage{} package~\cite{cgal:pt-t3-21a} provides several class templates, instances of which can be used to represent a variety 3D triangulations. In particular, the package provides the following class templates:
\begin{compactenum}
\item \cppLstinline{Triangulation\_3<Traits, Tds, Slds>},
\item \cppLstinline{Regular\_triangulation\_3<Traits, Tds, Slds>}, and
\item \cppLstinline{Delaunay\_triangulation\_3<Traits, Tds, LocationPolicy, Slds>}
\item \cppLstinline{Triangulation\_hierarchy\_3<Triangulation\_3>}
\end{compactenum}
Instance of the \cppLstinline{Delaunay\_triangulation\_3<>} and \cppLstinline{Regular\_triangulation\_3<>} class templates can be used to represent Delaunay and regular triangulation, respectively. In a regular triangulation points have an associated weight, and some points can be hidden and do not result in vertices in the triangulation. When instantiated its template parameter must be substituted with an instance of any other triangulation class template.

The \iiiDPeriodicTriangulationsPackage{} package~\cite{cgal:ct-pt3-21a} supports triangulations of sets of points in the three-dimensional flat torus~\cite{ct-cpt-09}. This package provides the class templates
\begin{compactenum}
\item \cppLstinline{Periodic\_3\_triangulation\_3<Traits, Tds>},
\item \cppLstinline{Periodic\_3\_regular\_triangulation\_3<Traits, Tds>}, and
\item \cppLstinline{Periodic\_3\_Delaunay\_triangulation\_3<Traits, Tds>}
\item \cppLstinline{Periodic\_3\_triangulation\_hierarchy\_3<PeriodicTriangulation>}
\end{compactenum}

The \cmake{} flag \myLstinline{TRI3\_NAME} specifies the particular types of triangulations, bindings for which should be generated; see Table~\ref{tab:tri3:names}.

\begin{table}[!htp]
  \caption{3D triangulation name options}
  \centering\begin{tikzpicture}
  \matrix (first) [tableName,
    column 1/.style={nodes={text width=10em}},
    column 2/.style={nodes={text width=39em}}
  ] {
    \myLstinline{TRI3\_NAME} & Type\\
    \myLstinline{plain} & \cppLstinline{Triangulation\_3<Traits, Tds, Slds>}\\
    \myLstinline{regular} & \cppLstinline{Regular\_triangulation\_3<Traits, Tds, Slds>}\\
    \myLstinline{delaunay} & \cppLstinline{Delaunay\_triangulation\_3<Traits, Tds, LocationPolicy, Slds>}\\
    \myLstinline{periodicPlain} & \cppLstinline{Periodic\_3\_triangulation\_3<Traits, Tds>}\\
    \myLstinline{periodicRegular} & \cppLstinline{Periodic\_3\_regular\_triangulation\_3<Traits, Tds>}\\
    \myLstinline{periodicDelaunay} & \cppLstinline{Periodic\_3\_Delaunay\_triangulation\_3<Traits, Tds>}\\
  };
  \end{tikzpicture}
  \label{tab:tri3:names}
\end{table}

%% \cppLstinline{TriangulationTraits\_3}
%% \cppLstinline{RegularTriangulationTraits\_3}
%% \cppLstinline{DelaunayTriangulationTraits\_3}
%% \cppLstinline{AlphaShapeTraits\_3}
%% \cppLstinline{FixedAlphaShapeTraits\_3}
%% \cppLstinline{WeightedAlphaShapeTraits\_3}
%% \cppLstinline{FixedWeightedAlphaShapeTraits\_3}
%% \cppLstinline{Periodic\_3TriangulationTraits\_3}
%% \cppLstinline{Periodic\_3RegularTriangulationTraits\_3}
%% \cppLstinline{Periodic\_3DelaunayTriangulationTraits\_3}
%% \cppLstinline{Periodic\_3AlphaShapeTraits\_3}
When any template above is instantiated the template parameter \cppLstinline{Traits} must be substituted with a model of a suitable geometric traits concept; this model is referred to as the geometric traits class; it provides the type of points to use as well as elementary operations on them. The type of traits used for the generated bindings is determined based on the selection of the triangulation type as explained below. It is conveniently defined in C++ as \cppLstinline{Dt::Geom\_traits}, where \cppLstinline{Dt} is the triangulation instance. This type is exposed as a Python attribute with the same name under \pyLstinline{Triangulation\_3} (which in turn is nested under the Python namespace \pyLstinline{Tri3}). Figure~\ref{tab:tri3:concept-hierarchy} depicts the 2D triangulation traits concept hierarchy. Any kernel instance is a model of any non-periodic traits concept; thus, when one of the templates
\begin{compactenum}
\item \cppLstinline{Triangulation\_3<Traits, Tds>},
\item \cppLstinline{Regular\_triangulation\_3<Traits, Tds>}, abd
\item \cppLstinline{Delaunay\_triangulation\_3<Traits, Tds>}
\end{compactenum}
is instantiated, the \cppLstinline{Traits} parameter is substituted with the \cppLstinline{Kernel} type (the selected kernel; see Section~\ref{ssec:modules:kernel}). When one of the templates
\begin{compactenum}
\item \cppLstinline{Periodic\_3\_triangulation\_3<Traits, Tds>},
\item \cppLstinline{Periodic\_3\_regular\_triangulation\_3<Traits, Tds>}, or
\item \cppLstinline{Periodic\_3\_Delaunay\_triangulation\_3<Traits, Tds>}
\end{compactenum}
is instantiated, the \cppLstinline{Traits} parameter is substituted
with \cppLstinline{Periodic\_3\_triangulation\_traits\_3<Kernel>},
\cppLstinline{Periodic\_3\_regular\_triangulation\_traits_3<Kernel>}, or
\cppLstinline{Periodic\_3\_Delaunay\_triangulation\_traits_3<Kernel>},
respectively.
%
Observe that when a Delaunay triangulation instance is used to define either a plain or a fixed alpha shape type (see Section~\ref{ssec:modules:as3}), the traits parameter must be substituted with a traits class that models either the \cppLstinline{AlphaShapeTraits\_3} or the
\cppLstinline{FixedAlphaShapeTraits\_3} concept, respectively. Similarly, when a regular triangulation instance is used to define either a plain or a fixed alpha shape type (see
Section~\ref{ssec:modules:as3}, the traits parameter must be substituted with a traits class that models either the \cppLstinline{WeightedAlphaShapeTraits\_3} or the \cppLstinline{FixedWeightedAlphaShapeTraits\_3} concept, respectively. Also in these cases the selected kernel serves as the traits.

\begin{figure}[!htp]
  \tikzset{
    my node style/.style={font=\scriptsize,minimum size=6mm,align=center,
      top color=white, bottom color=blue!25, draw=blue!75, drop shadow,
      rectangle, rounded corners, very thick
    }
  }
  \forestset{
    my tree style/.style={
      for tree={my node style,
        parent anchor=south,
        child anchor=north,
        edge={->,>=stealth,black,thick},
        edge path={% |-| style branches from parent to child
          \noexpand\path [draw, \forestoption{edge}] (!u.parent anchor)--(.child anchor)\forestoption{edge label};
        },
        if n children=3{% if a node has 3 children
          for children={% align the middle child with its parent
            if n=2{calign with current}{}
          }
        }{},
      }
    }
  }
  \centering
  \begin{forest}
    my tree style
    [ {\cppLstinline{TriTr3}}
      [ {\cppLstinline{RegularTriTr3}},name=r
        [ {\cppLstinline{WeightedASTr3}} ]
        [ {\cppLstinline{FixedWeightedASTr3}} ]
      ]
      [ {\cppLstinline{Periodic\_3TriTr3}},name=p
        [ {\cppLstinline{Pr3RegularTriTr3}},name=pr ]
      ]
      [ {\cppLstinline{DelaunayTriTr3}},name=d
        [ {\cppLstinline{Pr3DelaunayTriTr3}},name=pd
	  [ {\cppLstinline{Pr3ASTr3}} ]
        ]
        [ {\cppLstinline{ASTr3}} ]
        [ {\cppLstinline{FixedASTr3}} ]
      ]
    ]
    \draw[->,>=stealth] (r.parent anchor)--(pr.child anchor);
    \draw[->,>=stealth] (p.parent anchor)--(pd.child anchor);
  \end{forest}
  \label{tab:tri3:concept-hierarchy}
  \caption{The 3D triangulation traits concept hierarchy.
    \cppLstinline{Triangulation}, \cppLstinline{AlphaShape},
    \cppLstinline{Traits\_3} and \cppLstinline{Periodic\_3} are
    abbreviated as \cppLstinline{Tri}, \cppLstinline{AS},
    \cppLstinline{Pr3}, and \cppLstinline{Tr3}, respectively.}
\end{figure}

The type that substitutes the \cppLstinline{Tds} parameter models the concept
\cppLstinline{TriangulationDataStructure\_3}. An object of this type stores the combinatorial structure of the triangulation; it is an instance of the class template \cppLstinline{Triangulation\_data\_structure\_3<VertexBase, CellBase, ConcurrencyTag>}.

The types that substitute the \cppLstinline{VertexBase} and \cppLstinline{CellBase} template parameters when the template \cppLstinline{Triangulation\_data\_structure\_3<VertexBase, CellBase>} is instantiated represent the type of the vertex and the type of the cell of the triangulation, respectively; they must model the concepts \cppLstinline{TriangulationDSVertexBase\_3} and \cppLstinline{TriangulationDSCellBase\_3}, respectively. If the binding is generated for a periodic triangulation, these parameters must be substituted with types that also model the concepts \cppLstinline{Periodic\_3TriangulationDSVertexBase\_3} and \cppLstinline{Periodic\_3TriangulationDSCellBase\_3}, respectively. If the triangulation is used to define a plain alpha shape type, these parameters must be substituted with types that model the concepts \cppLstinline{AlphaShapeVertex\_3} and \cppLstinline{AlphaShapeCell\_3}, respectively; see Section~\ref{ssec:modules:as3}. If the triangulation is used to define a fixed alpha shape type, these parameters must be substituted with types that model the concepts \cppLstinline{FixedAlphaShapeVertex\_3} and \cppLstinline{FixedAlphaShapeCell\_3}, respectively. Finally, if the binding is generated for a hierarchy triangulation, e.g., an instance of the template parameter \cppLstinline{Triangulation\_hierarchy\_3<Triangulation\_3>}, the \cppLstinline{VertexBase} template parameter must be substituted with a model of the concept
\cppLstinline{TriangulationHierarchyVertexBase\_3}. (Observe that there is no special requirements on the type that substitutes the \cppLstinline{CellBase} parameter in this case.) Below we describe how the final vertex type is selected in details.

%% 3D Triangulation Vertex Selection
If bindings is generated for a periodic triangulation, the type \cppLstinline{Vbp} is defined as
\cppLstinline{Periodic\_3\_triangulation\_ds\_vertex\_base\_2<>}; otherwise, \cppLstinline{Vbp} is defined as \cppLstinline{Triangulation\_ds\_vertex\_base\_3<>}.
%
If bindings is generated for an instance of either \cppLstinline{Regular\_triangulation\_vertex\_base\_3<Kernel>} or \cppLstinline{Periodic\_3_Regular\_triangulation\_vertex\_base\_3<Kernel>}, the type \cppLstinline{Vb} is defined as \cppLstinline{Regular\_triangulation\_vertex\_base\_2<Traits, Vbp>}; otherwise, the type \cppLstinline{Vb} is defined as \cppLstinline{Triangulation\_vertex\_base\_3<Traits, Vbp>}.
%
If the \cmake{} flag \cmakeLstinline{TRI2\_VERTEX\_WITH\_INFO} is set to \myLstinline{true}, the type \cppLstinline{Vbi} is defined as \cppLstinline{Triangulation\_vertex\_base\_with\_info\_3<Data, Traits, Vb>}, where \cppLstinline{Data} is defined as the generic Python object \cppLstinline{bp::object}; see Section~\ref{sec:bind:ae}; otherwise the type \cppLstinline{Vbi} is defined as \cppLstinline{Vb}.
%
If the \cmake{} flag \cmakeLstinline{TRI2\_HIERARCHY} is set to \myLstinline{true}, the type \cppLstinline{Vbih} is defined as \cppLstinline{Triangulation\_hierarchy\_vertex\_base\_2<Vbi>}; otherwise the type \cppLstinline{Vbih} is defined as \cppLstinline{Vbi}.
%
Finally, if the triangulation is used to define an alpha shape type, the final type \cppLstinline{V} (that substitutes the \cppLstinline{VertexBase} parameter) is defined as either
\cppLstinline{Alpha\_shape\_vertex\_base\_2<Traits, Vbih, ExactComparison>} or
\cppLstinline{Fixed\_Alpha\_shape\_vertex\_base\_2<Traits, Vbih, ExactComparison>} based on the selection of the setting of the \cmake{} flag \cmakeLstinline{AS3\_NAME}; see Table~\ref{tab:as3};
\cppLstinline{ExactComparison} is substituted with either \cppLstinline{true} or \cppLstinline{false} based on the setting of the \cmake{} flag \cmakeLstinline{AS3\_EXACT\_COMPARISON}; see Table~\ref{tab:as2}. Otherwise, \cppLstinline{V} is defined as \cppLstinline{Vbih}.
%
The final vertex type is conveniently defined in C++ as \cppLstinline{Dt::Vertex}, where \cppLstinline{Dt} is the triangulation instance. This type is exposed as a Python attribute
with the same name under \pyLstinline{Triangulation_2}.

The final face type is determined in a similar fashion. The only difference between the selections of the vertex and face types is that the setting of the \cmake{} flag \cmakeLstinline{TRI2\_HIERARCHY} does not have an effect on the face type selection. The final face type is conveniently defined in C++ as \cppLstinline{Dt::Face}, where \cppLstinline{Dt} is the triangulation instance. This type is also exposed as a Python attribute with the same name under \pyLstinline{Triangulation_2}.

The \cmake{} Boolean flag \cmakeLstinline{TRI2\_HIERARCHY} indicates whether the type of triangulation, bindings for which should be generated, is an instance of one of the triangulation hierarchy templates, namely, \cppLstinline{Triangulation\_hierarchy\_2<Triangulation\_2>} and
\cppLstinline{Periodic_triangulation\_hierarchy\_2<PeriodicTriangulation\_2>}. The template parameter in both cases is substituted with the triangulation selected via the \cmake{} flag
\cmakeLstinline{TRI2\_NAME}, which also determines whether to use the non-periodic or periodic version above.

The template parameter \cppLstinline{ConcurrencyTag} is substituted with either \cppLstinline{Sequential\_tag} or \cppLstinline{Parallel\_tag} when the template \cppLstinline{Triangulation\_data\_structure\_3<VertexBase, CellBase, ConcurrencyTag} is instantiated. It enables the use of a concurrent container to store vertices and cells. The \cmake{} flag \cmakeLstinline{TRI3\_CONCURRENCY\_NAME} determines the selection; see Table~\ref{tab:tri3:concurrency}.

\begin{table}[!htp]
  \caption{3D triangulation concurrency options}
  \centering\begin{tikzpicture}
  \matrix (first) [tableName,
    column 1/.style={nodes={text width=18em}},
    column 2/.style={nodes={text width=31em}}
  ] {
    \cmakeLstinline{TRI3\_CONCURRENCY\_NAME} & Type\\
    \myLstinline{sequential} & \cppLstinline{Sequential\_tag}\\
    \myLstinline{parallel} & \cppLstinline{Parallel\_tag}\\
  };
  \end{tikzpicture}
  \label{tab:tri3:concurrency}
\end{table}

The template parameter \cppLstinline{LocationPolicy} is substituted with either \cppLstinline{Fast\_location} or \cppLstinline{Compact\_location} when the template
\cppLstinline{Delaunay\_triangulation\_3<Traits, Tds, LocationPolicy>} is instantiated. It enables a faster point location at the account of memory space. This is useful when performing point locations or random point insertions in large data sets.  The \cmake{} flag \cmakeLstinline{TRI3\_LOCATION\_POLICY\_NAME} determines the selection; see Table~\ref{tab:tri3:location}.

\begin{table}[!htp]
  \caption{3D triangulation location policy options}
  \centering\begin{tikzpicture} \matrix (first) [tableName, column
    1/.style={nodes={text width=18em}}, column 2/.style={nodes={text
        width=31em}} ] { \cmakeLstinline{TRI3\_LOCATION\_POLICY\_NAME}
    & Type\\ \myLstinline{fast} &
    \cppLstinline{Fast\_location}\\ \myLstinline{compact} &
    \cppLstinline{Compact\_location}\\ };
  \end{tikzpicture}
  \label{tab:tri3:location}
\end{table}

The following code excerpt in the \cmake{} language sets the \cmake{} flags for our next example.
\framedInputCmake[\normalsize]{cmake/tri3_del_epic_release.cmake}

The following example constructs a three-dimensional Delaunay triangulation from six points and verifies that the triangulation is valid.
\framedInputPython[\normalsize][15]{python/tri3_del.py}

% -----------------------------------------------------------------------------
\subsection{3D Alpha Shape Bindings}
\label{ssec:modules:as3}
% -----------------------------------------------------------------------------
Given a set of points sampled in a 3D body, an alpha shape (a.k.a.\ alpha complex) is demarcated by a frontier, which is a linear approximation of the original boundary of the body. Similar to the 2D alpha shape object, A 3D alpha shape object maintains an underlying 3D triangulation of a set of input points. Basic alpha shapes are based on the Delaunay triangulation and weighted alpha shapes are based on regular triangulation, where the euclidean distance is replaced by the
power to weighted points. The package \iiiDAlphaShapesPackage{}~\cite{cgal:dy-as3-21a} provides the class templates \cppLstinline{Alpha\_shape\_3<Dt, ExactAlphaComparisonTag>} and \cppLstinline{Fixed\_alpha\_shape\_3<Dt>}. Instances of both templates can be used to represent a large variety of alpha shapes for a given set of points. In a plain alpha shape each $k$-face of this triangulation is associated with an interval specifying for which values of $\alpha$ the face belongs to the alpha complex. In a fixed alpha shape each $k$-face is associated with a classification that specifies its status in the alpha complex, alpha being fixed. Table~\ref{tab:as3} lists the \cmake{} flags associated with the \iiiDAlphaShapesPackage{} module.

%%%%%%%% 3D Alpha Shape flags
\begin{table}[!htp]
  \caption{\iiiDAlphaShapesPackage{} module flags}
  \centering\begin{tikzpicture}
  \matrix (first) [tableModuleFlags,
    column 1/.style={nodes={text width=16em}},
    column 3/.style={nodes={text width=5.7em}},
    column 4/.style={nodes={text width=21.6em}}
  ] {
Name & Type &  Default & Description\\
\myLstinline{AS3\_NAME}              & String  & \myLstinline{plain} & The 3D Alpha shape type\\
\myLstinline{AS3\_EXACT\_COMPARISON} & Boolean & \myLstinline{false} & Determines whether to apply exact comparisons\\
  };
 \end{tikzpicture}
 \label{tab:as3}
\end{table}

The \cmake{} flag \myLstinline{AS3\_NAME} specifies which alpha shape
should be used for the generated bindings; see
Table~\ref{tab:as3:types}

%%%%%%%% 3D Alpha Shapes Module Types
\begin{table}[!htp]
 \caption{\iiiDAlphaShapesPackage{} types}
  \centering\begin{tikzpicture}
  \matrix (first) [tableName,
    column 1/.style={nodes={text width=12em}},
    column 2/.style={nodes={text width=37em}}
  ] {
     \myLstinline{AS3\_NAME} & Type\\
     \myLstinline{plain} & \cppLstinline{Alpha\_shape\_3<>}\\
     \myLstinline{fixed} & \cppLstinline{Fixed\_alpha\_shape\_3<>}\\
  };
 \end{tikzpicture}
 \label{tab:as3:types}
\end{table}

The template parameter \cppLstinline{Dt} must be substituted with an
instance of one of the following class templates that can be used to
represent a triangulation:
\begin{compactenum}
\item \cppLstinline{Delaunay\_triangulation\_3},
\item \cppLstinline{Regular\_triangulation\_3},
\item \cppLstinline{Periodic\_3_Delaunay_triangulation\_3}, or
\item \cppLstinline{Periodic_3\_regular_triangulation_3};
\end{compactenum}
see Section~\ref{ssec:modules:tri3}.
Note that \cppLstinline{Dt::Geom\_traits} must be model of a suitable
alpha-shape traits concept; see Table~\ref{tab:tri3:traits}, and
\cppLstinline{Dt::Vertex} and \cppLstinline{Dt::Face} must be models
suitable alpha-shape vertex and cell concepts, respectively; see
Table~\ref{tab:tri3:entity}.

The following code excerpt in the \cmake{} language sets the \cmake{}
flags for our next example.
\framedInputCmake[\normalsize]{cmake/as3_del_compact_parallel_epic_release.cmake}

The following example constructs an alpha shape object.
\framedInputPython[\normalsize][15]{python/as3_del_compact_epic.py}

% -----------------------------------------------------------------------------
\subsection{Spatial Searching Bindings}
\label{ssec:modules:ss}
% -----------------------------------------------------------------------------
The \dDSpatialSearchingPackage{} package~\cite{cgal:tf-ssd-21a} implements exact and approximate distance browsing by providing implementations of algorithms supporting
\begin{compactenum}
\item both nearest and furthest neighbor searching,
\item both exact and approximate searching,
\item (approximate) range searching,
\item (approximate) k-nearest and k-furthest neighbor searching,
\item (approximate) incremental nearest and incremental furthest neighbor searching, and
\item query items representing points and spatial objects.
\end{compactenum}
In these searching problems a set $\calP$ of data points in $d$-dimensional space is given. The points in $\calP$ are preprocessed into a tree data structure, so that given any query item $q$ the points of $\calP$ can be browsed efficiently. The approximate \dDSpatialSearchingPackage{} package is designed for data sets that are small enough to store the search structure in main memory (in contrast to approaches from databases that assume that the data reside in secondary storage).

%%%%%%%%
\begin{table}[!htp]
  \caption{\dDSpatialSearchingPackage{} module flags}
  \centering\begin{tikzpicture}
  \matrix (first) [tableModuleFlags,
    column 3/.style={nodes={text width=3.5em}},
    column 4/.style={nodes={text width=21.8em}}
  ] {
    Name & Type &  Default & Description\\
    \cmakeLstinline{SPATIAL\_SEARCHING\_DIMENSION} & Integer & \cppLstinline{2} &
      The dimension of spatial searching related classes\\
  };
 \end{tikzpicture}
 \label{tab:ss}
\end{table}

\framedInputCmake[\normalsize]{cmake/ss_2_cdlg_release.cmake}

\framedInputPython[\normalsize][15][67]{python/kd_tree.py}
%% 2. kernel-d is used by spatial searching
